\documentclass[11pt, a4paper]{article}

% bfw-style.tex — gemeinsame Style-Definitionen für alle Handouts/Aufgabenhefte
% Voraussetzung: Kompilation mit xelatex (wegen fontspec)

\usepackage[ngerman]{babel}
\usepackage{geometry}
\geometry{a4paper, top=2.2cm, bottom=2.2cm, left=2.3cm, right=2.3cm}

\usepackage{fontspec}
% Fallback-Kette: Source Sans Pro -> Calibri -> Segoe UI -> Latin Modern Sans
% (Latin Modern ist immer mit MiKTeX vorhanden)
\IfFontExistsTF{Source Sans Pro}{%
  \setmainfont{Source Sans Pro}[Ligatures=TeX]%
}{\IfFontExistsTF{Calibri}{%
  \setmainfont{Calibri}[Ligatures=TeX]%
}{\IfFontExistsTF{Segoe UI}{%
  \setmainfont{Segoe UI}[Ligatures=TeX]%
}{%
  \setmainfont{Latin Modern Sans}[Ligatures=TeX]%
}}}
\IfFontExistsTF{Source Code Pro}{%
  \setmonofont{Source Code Pro}[Scale=0.92]%
}{\IfFontExistsTF{Consolas}{%
  \setmonofont{Consolas}[Scale=0.92]%
}{%
  \setmonofont{Latin Modern Mono}[Scale=0.92]%
}}

\usepackage{xcolor}
% Farbpalette: hell, freundlich, modern
\definecolor{bfwPrimary}{HTML}{0EA5A4}    % Petrol/Türkis - Primär
\definecolor{bfwPrimaryDark}{HTML}{0F766E} % dunkler Petrol für Headings
\definecolor{bfwAccent}{HTML}{F59E0B}     % Amber - Akzent
\definecolor{bfwSuccess}{HTML}{10B981}    % Grün - Erfolg / Lösung
\definecolor{bfwWarn}{HTML}{F97316}       % Orange - Achtung
\definecolor{bfwInk}{HTML}{0F172A}        % fast Schwarz
\definecolor{bfwInkSoft}{HTML}{475569}    % gedämpftes Grau
\definecolor{bfwBgSoft}{HTML}{F8FAFC}     % off-white
\definecolor{bfwBgTeal}{HTML}{ECFEFF}     % helles Türkis
\definecolor{bfwBgAmber}{HTML}{FEF3C7}    % helles Gelb
\definecolor{bfwBgRose}{HTML}{FFE4E6}     % helles Rosé für Eselsbrücken

\usepackage[colorlinks=true, linkcolor=bfwPrimaryDark, urlcolor=bfwPrimaryDark]{hyperref}
\usepackage{enumitem}
\setlist[itemize]{leftmargin=*, itemsep=2pt, topsep=2pt}
\setlist[enumerate]{leftmargin=*, itemsep=2pt, topsep=2pt}

\usepackage[most]{tcolorbox}
\usepackage{booktabs}
\usepackage{tikz}
\usetikzlibrary{decorations.pathmorphing, positioning, shapes.geometric, arrows.meta}

% Merkkästchen — wichtige Definition / Kernpunkt
\newtcolorbox{merksatz}[1][]{%
  enhanced, breakable,
  colback=bfwBgTeal,
  colframe=bfwPrimary,
  boxrule=0.6pt,
  arc=4pt,
  left=10pt, right=10pt, top=8pt, bottom=8pt,
  fonttitle=\bfseries\color{bfwPrimaryDark},
  title={\faLightbulb\ Merksatz},
  #1
}

% Eselsbrücke — Mnemoniks
\newtcolorbox{eselsbruecke}[1][]{%
  enhanced, breakable,
  colback=bfwBgRose,
  colframe=bfwAccent,
  boxrule=0.6pt,
  arc=4pt,
  left=10pt, right=10pt, top=8pt, bottom=8pt,
  fonttitle=\bfseries\color{bfwWarn},
  title={Eselsbrücke},
  #1
}

% Beispiel-Box
\newtcolorbox{beispiel}[1][]{%
  enhanced, breakable,
  colback=bfwBgSoft,
  colframe=bfwInkSoft,
  boxrule=0.4pt,
  arc=3pt,
  left=10pt, right=10pt, top=8pt, bottom=8pt,
  fonttitle=\bfseries\color{bfwInk},
  title={Beispiel},
  #1
}

% Lösungs-Box (für Aufgabenhefte)
\newtcolorbox{loesung}[1][]{%
  enhanced, breakable,
  colback=bfwBgAmber!40,
  colframe=bfwSuccess,
  boxrule=0.6pt,
  arc=3pt,
  left=10pt, right=10pt, top=8pt, bottom=8pt,
  fonttitle=\bfseries\color{bfwSuccess!70!black},
  title={Lösung},
  #1
}

% Glossar-Eintrag
\newcommand{\glossareintrag}[2]{%
  \par\noindent\textbf{\color{bfwPrimaryDark}#1} \enspace #2\par\smallskip
}

% Section-Stil — etwas farbiger als Standard
\usepackage{titlesec}
\titleformat{\section}
  {\Large\bfseries\color{bfwPrimaryDark}}
  {\thesection}{0.6em}{}
\titleformat{\subsection}
  {\large\bfseries\color{bfwPrimary}}
  {\thesubsection}{0.5em}{}
\titleformat{\subsubsection}
  {\normalsize\bfseries\color{bfwInk}}
  {\thesubsubsection}{0.4em}{}

% TikZ-Stil "skizzenhaft" — etwas weicher als Rough.js, aber stabil
\tikzset{
  roughbox/.style = {
    draw=bfwPrimaryDark, thick, fill=bfwBgTeal,
    rounded corners=4pt, inner sep=6pt
  },
  roughaccent/.style = {
    draw=bfwAccent, thick, fill=bfwBgAmber,
    rounded corners=4pt, inner sep=6pt
  },
  roughline/.style = {
    draw=bfwInkSoft, thick, -{Stealth[length=2.5mm]}
  }
}

% Bullet-Symbole (bewusst kein FontAwesome — vermeidet Font-Installations-Probleme)
\newcommand{\faLightbulb}{$\bullet$}
\newcommand{\faBookmark}{$\bullet$}

% Header/Footer
\usepackage{fancyhdr}
\pagestyle{fancy}
\fancyhf{}
\renewcommand{\headrulewidth}{0pt}
\fancyfoot[L]{\small\color{bfwInkSoft}\thepage}
\fancyfoot[R]{\small\color{bfwInkSoft} BFW M-064 Beschaffung im Büromanagement}


\title{\vspace*{-1.5cm}\Huge\bfseries\color{bfwPrimaryDark}Tag~5: Annahme und Lagerung der Waren\\[0.4em]
  \large\color{bfwInkSoft}\textnormal{Wareneingangskontrolle, § 377 HGB, Lagerformen, Lagerkennzahlen}}
\author{\textsf{Beschaffung im Büromanagement}\\
\textsf{\small\copyright\ Can Siebert\,\textperiodcentered\,2026}}
\date{18.05.2026}

\begin{document}
\maketitle
\thispagestyle{empty}

\vspace{1em}
\begin{merksatz}[title={\bfseries Worum geht es heute?}]
Die Lieferung kommt --- was nun? Heute gehen wir den Wareneingang strukturiert durch:
Wie nehmen Sie die Ware an, was prüfen Sie wann, und wie reagieren Sie bei Mängeln?
Anschließend wenden wir uns der Lagerhaltung zu --- ihren Funktionen, Formen, Arten,
Kosten und vor allem den \emph{Lagerkennzahlen}, die in fast jeder IHK-Klausur Teil~2
abgefragt werden.
\end{merksatz}

\tableofcontents
\vspace{1em}
\hrule
\vspace{1em}

\section{Wareneingangskontrolle}

\subsection{Der Ablauf in vier Schritten}

Sobald die Lieferung am Wareneingang ankommt, beginnt ein klar strukturierter Prozess:

\textbf{Schritt 1 --- Annahme}: Lieferung entgegennehmen. Lieferschein unterschreiben ---
allerdings nur, wenn Sie keine offensichtlichen Probleme sehen. Bei sichtbaren Schäden:
\textbf{mit Vorbehalt} unterschreiben (z.\,B.\ ,,unter Vorbehalt der Mängelprüfung'').

\textbf{Schritt 2 --- Äußere Kontrolle}: Vor dem Fahrer Anzahl der Pakete zählen,
sichtbare Beschädigungen am Karton oder der Folie protokollieren. Diese Kontrolle
sichert Sie gegen Transportschäden ab --- der Fahrer ist Zeuge.

\textbf{Schritt 3 --- Innere Kontrolle}: Nach Abfahrt des Fahrers auspacken. Stichprobe
oder vollständige Prüfung, je nach Wert und Kritikalität. Vergleich mit Lieferschein
und Bestellung. Gefunden werden Mängel hinsichtlich \emph{Menge}, \emph{Beschaffenheit}
und \emph{Identität} der Ware.

\textbf{Schritt 4 --- Erfassung}: Erfassung im Warenwirtschafts­system, Einlagerung am
bestimmten Lagerplatz. Erst jetzt ist der Wareneingang abgeschlossen.

\subsection{Offene und verdeckte Mängel}

Mängel werden rechtlich in zwei Kategorien eingeteilt:

\emph{Offene Mängel} sind bei der üblichen Wareneingangs­prüfung sofort erkennbar.
Beispiele: beschädigtes Verpackungs­material, falsche Stückzahl, sichtbare Verformungen,
falsche Farbe.

\emph{Verdeckte Mängel} zeigen sich erst später, oft im Gebrauch. Beispiele: defekte
Mechanik in einem Aktenordner, schlechter Druck auf vermeintlich gutem Druckerpapier,
Feuchtigkeitsschäden, die erst beim Auspacken einer hinteren Lage entdeckt werden.

\textbf{Beide Mängelarten} müssen gerügt werden --- offene unverzüglich nach Lieferung,
verdeckte unverzüglich nach Entdeckung.

\section{§ 377 HGB --- die Rügepflicht}

\subsection{Was das Gesetz verlangt}

\begin{merksatz}[title={§ 377 HGB im Wortlaut (sinngemäß)}]
Beim Handelskauf (Kauf zwischen Kaufleuten) muss der Käufer die Ware \emph{unverzüglich}
nach Ablieferung untersuchen, soweit dies nach ordnungsmäßigem Geschäftsgang tunlich ist.
Zeigt sich ein Mangel, ist der Verkäufer \emph{unverzüglich} zu benachrichtigen.
\end{merksatz}

Das ist eine sogenannte \emph{Obliegenheit} --- keine Pflicht im engeren Sinn, sondern
eine Last: Wer sie versäumt, verliert seine Rechte.

\subsection{Folge bei Versäumnis}

Versäumt der Käufer die rechtzeitige Mängelrüge, gilt die Ware nach §\,377 Abs.\,2 HGB
als \emph{genehmigt}. Konkret bedeutet das: Der Käufer verliert sämtliche Mängelrechte,
also Nacherfüllung, Rücktritt, Minderung und Schadensersatz.

\begin{eselsbruecke}
\textbf{,,Wer schweigt, hat zugestimmt.''} Das gilt im Handelsverkehr besonders streng.
Im B2B muss man bei Mängeln aktiv werden --- still sein bedeutet Verlust der Rechte.
\end{eselsbruecke}

\subsection{Was ist „unverzüglich"?}

§\,121 BGB definiert ,,unverzüglich'' als ,,ohne schuldhaftes Zögern''. Die Rechtsprechung
hat das in der Praxis meist auf 1--2~Werktage konkretisiert. Bei komplizierter Prüfung
(z.\,B.\ technische Geräte) kann auch eine Woche akzeptabel sein.

\subsection{Inhalt einer Mängelrüge}

Eine vollständige Mängelrüge enthält:

\begin{itemize}[itemsep=0pt]
  \item Bezug auf die Lieferung (Lieferschein-Nr., Datum, Bestell-Nr.)
  \item Genaue Beschreibung des Mangels
  \item Konkretes Gewähr\-leistungs\-recht, das ausgeübt wird (Nacherfüllung, Minderung, Rücktritt, Schadensersatz)
  \item Frist für die Reaktion des Verkäufers
  \item Datum und Unterschrift
\end{itemize}

Form: schriftlich (E-Mail, Brief, Fax). Ein bloßer Telefonanruf reicht aus
Beweisgründen nicht.

\section{Aufgaben und Funktionen der Lagerhaltung}

Warum überhaupt Lager halten? Die Lagerhaltung erfüllt vier zentrale Funktionen:

\begin{itemize}[itemsep=0pt]
  \item \textbf{Sicherungsfunktion}: Schutz gegen Lieferunterbrechungen, Streiks, Bedarfsspitzen.
  \item \textbf{Bereitstellungsfunktion}: Material steht für die Produktion oder den Verkauf bereit, ohne auf jeden Lieferanten warten zu müssen.
  \item \textbf{Überbrückungsfunktion}: Zwischenlagerung bei zeitlicher Differenz zwischen Beschaffung und Verbrauch (z.\,B.\ Einkauf in größeren Mengen, sukzessiver Verbrauch).
  \item \textbf{Spekulationsfunktion}: Einkauf zu günstigem Zeitpunkt (z.\,B.\ vor erwarteter Preiserhöhung) und Lagerung für späteren Verbrauch.
\end{itemize}

Daneben gibt es die \emph{Veredelungsfunktion} bei bestimmten Materialien, die durch
Lagerung ihren Wert oder ihre Eigenschaften verbessern (Wein, Käse, Holz).

\section{Lagerformen und Lagerarten}

\subsection{Eigenlager oder Fremdlager?}

\begin{center}
\begin{tabular}{p{4cm} p{5cm} p{5cm}}
\toprule
\textbf{Kriterium} & \textbf{Eigenlager} & \textbf{Fremdlager}\\
\midrule
Investition & Hoch (Gebäude, Einrichtung) & Keine\\
Fixkosten & Hoch & Variable Kosten\\
Flexibilität & Gering (Größe steht fest) & Hoch (Anpassung an Bedarf)\\
Kontrolle & Volle Kontrolle & Eingeschränkt\\
Know-how & Beim eigenen Personal & Beim Dienstleister\\
Geeignet für & Großvolumiger Dauerbedarf & Schwankender oder kleiner Bedarf\\
\bottomrule
\end{tabular}
\end{center}

In der Praxis wählen viele Unternehmen Mischformen: Stammbestände im Eigenlager,
Saisonspitzen im Fremdlager.

\subsection{Lagerarten nach Bauart}

\begin{itemize}[itemsep=0pt]
  \item \textbf{Bodenlager}: Waren stehen direkt am Boden. Einfach, günstig, geeignet für Großgut.
  \item \textbf{Regallager}: Standardregale, manuell oder mit Stapler. Flexibel, klassisch im Mittelstand.
  \item \textbf{Hochregallager}: Bis zu 50~m hoch, automatisiert. Maximale Raumnutzung, hohe Investition. Typisch für große Distributionszentren.
  \item \textbf{Blocklager}: Stapelung von Paletten ohne Regal. Günstig, geeignet für gleichartige Waren in großer Menge.
  \item \textbf{Kanallager}: Paletten in Kanälen, Entnahme nach FIFO oder LIFO. Hohe Flächennutzung.
  \item \textbf{Etagenlager}: Mehrere Stockwerke nutzen --- für innerstädtische Lager mit knappen Grundstücken.
\end{itemize}

\section{Lagerkosten}

Lagerhaltung ist nie umsonst. Die Lagerkosten setzen sich aus mehreren Komponenten
zusammen:

\begin{itemize}[itemsep=0pt]
  \item \emph{Personalkosten}: Lagermitarbeiter, Schichtleitung, Disposition.
  \item \emph{Raumkosten}: Lagermiete, Energie, Reinigung.
  \item \emph{Anlagekosten}: Abschreibung von Regalen, Staplern, IT-Systemen.
  \item \emph{Kapitalbindungskosten}: Entgangene Zinsen für das Geld, das in Lagerbeständen gebunden ist (Opportunitätskosten).
  \item \emph{Versicherungskosten}: Sachversicherungen für den Lagerwert.
  \item \emph{Lagerrisikokosten}: Verderb, Schwund, Diebstahl, technisches Veralten.
\end{itemize}

\begin{merksatz}
Die \textbf{Kapitalbindung} ist oft der größte versteckte Kostenblock. Geld im Lager
liegt brach --- es kann nicht in andere Investitionen fließen oder Zinsen auf der Bank
generieren. Diese Opportunitätskosten erfasst der \emph{Lagerzinssatz}.
\end{merksatz}

\section{Lagerkennzahlen --- die wichtigsten vier}

Vier Kennzahlen müssen Sie für die IHK-Prüfung sicher beherrschen.

\subsection{Durchschnittlicher Lagerbestand}

\begin{merksatz}[title={Formel}]
\textbf{Ø Lagerbestand = (Anfangsbestand + Endbestand) $\div$ 2}
\end{merksatz}

Genauere Variante mit 12 Monatsbeständen:

\[
\varnothing\,\text{Lagerbestand} = \frac{\text{AB} + \text{B}_1 + \text{B}_2 + \cdots + \text{B}_{12}}{13}
\]

\begin{beispiel}
Anfangsbestand 4\,000~Stück, Endbestand 6\,000~Stück. Ø Lagerbestand $= (4\,000 + 6\,000) \div 2 = \textbf{5\,000 Stück}$. Bei 5,00~€ pro Stück: \textbf{25\,000~€} Ø Lagerwert.
\end{beispiel}

\subsection{Lagerumschlagshäufigkeit (LUH)}

\begin{merksatz}[title={Formel}]
\textbf{LUH = Jahresverbrauch (oder Wareneinsatz) $\div$ Ø Lagerbestand}
\end{merksatz}

Sagt: Wie oft pro Jahr wird der Lagerbestand komplett umgeschlagen?

\begin{beispiel}
Jahresverbrauch 50\,000~Stück, Ø Lagerbestand 5\,000~Stück $\to$ LUH $= 50\,000 \div 5\,000 = \textbf{10}$. Das Lager wird also 10-mal pro Jahr umgeschlagen.
\end{beispiel}

\textbf{Interpretation}: Hohe LUH = wenig Kapitalbindung = wirtschaftlich gut. Niedrige
LUH = Lager liegt zu lange. Branchenwerte: Lebensmitteleinzelhandel oft LUH > 30,
Maschinenbau oft LUH 4--8.

\subsection{Durchschnittliche Lagerdauer}

\begin{merksatz}[title={Formel}]
\textbf{Ø Lagerdauer (in Tagen) = 360 $\div$ LUH}
\end{merksatz}

Sagt: Wie lange liegt eine Ware durchschnittlich im Lager? Hinweis: Das kaufmännische
Jahr hat 360~Tage (12~$\times$~30), nicht 365.

\begin{beispiel}
Bei LUH = 10 $\to$ Ø Lagerdauer $= 360 \div 10 = \textbf{36 Tage}$.
\end{beispiel}

\subsection{Lagerzinsen}

\begin{merksatz}[title={Formel (vereinfacht für Jahr)}]
\textbf{Lagerzinsen = Ø Lagerwert $\times$ Lagerzinssatz}
\end{merksatz}

\begin{beispiel}
Ø Lagerwert 25\,000~€, Lagerzinssatz 8\,\% $\to$ Lagerzinsen $= 25\,000 \times 0{,}08 = \textbf{2\,000 €} \text{ pro Jahr}$. Diese 2\,000~€ sind die ,,versteckten'' Kosten der Kapitalbindung.
\end{beispiel}

Genauere Formel (für unterjährig):
\[
\text{Lagerzinsen} = \frac{\text{Ø Lagerwert} \times \text{Lagerdauer} \times \text{Lagerzinssatz}}{360}
\]

\section{Zusammenfassung}

\begin{enumerate}
  \item \textbf{Wareneingang}: Annahme $\to$ äußere $\to$ innere Kontrolle $\to$ Erfassung.
  \item \textbf{§ 377 HGB}: Beim Handelskauf unverzüglich prüfen und rügen --- sonst Verlust der Mängelrechte.
  \item \textbf{Mängel}: offen oder verdeckt --- beide unverzüglich rügen.
  \item \textbf{Lagerformen}: Eigenlager (Kontrolle, hohe Fixkosten) vs.\ Fremdlager (flexibel, variable Kosten).
  \item \textbf{Lagerarten}: Boden, Regal, Hochregal, Block, Kanal, Etage.
  \item \textbf{Lagerkennzahlen}: Ø~Bestand $=$ (AB$+$EB)$\div$2; LUH $=$ Verbrauch $\div$ Ø~Bestand; Ø~Dauer $=$ 360$\div$LUH; Lagerzinsen $=$ Ø~Wert $\times$ Zinssatz.
\end{enumerate}

\newpage

\section*{Glossar}
\addcontentsline{toc}{section}{Glossar}

\glossareintrag{Wareneingangskontrolle}{Strukturierte Prüfung der angelieferten Ware in vier Schritten: Annahme, äußere Kontrolle, innere Kontrolle, Erfassung.}

\glossareintrag{Äußere Kontrolle}{Sichtprüfung der Verpackung und Anzahl der Pakete vor dem Frachtfahrer --- Schutz vor Transportschäden.}

\glossareintrag{Innere Kontrolle}{Prüfung des Inhalts (Menge, Beschaffenheit, Identität) nach Abfahrt des Fahrers --- Vergleich mit Lieferschein und Bestellung.}

\glossareintrag{Mängelrüge}{Schriftliche Anzeige eines festgestellten Mangels an den Verkäufer mit konkreter Beschreibung und gewähltem Gewährleistungsrecht.}

\glossareintrag{§ 377 HGB}{Rügeobliegenheit beim Handelskauf --- der Käufer muss unverzüglich prüfen und Mängel anzeigen, sonst verliert er seine Mängelrechte.}

\glossareintrag{Offener Mangel}{Mangel, der bei ordnungsgemäßer Wareneingangsprüfung sofort erkennbar ist.}

\glossareintrag{Verdeckter Mangel}{Mangel, der erst später (z.\,B.\ im Gebrauch) entdeckt wird --- muss unverzüglich nach Entdeckung gerügt werden.}

\glossareintrag{Sicherungsfunktion}{Aufgabe des Lagers, gegen Lieferunterbrechungen und Bedarfsspitzen abzusichern.}

\glossareintrag{Bereitstellungsfunktion}{Aufgabe des Lagers, Material bereitzuhalten, sodass es jederzeit verfügbar ist.}

\glossareintrag{Überbrückungsfunktion}{Aufgabe des Lagers, zeitliche Differenzen zwischen Beschaffung und Verbrauch zu überbrücken.}

\glossareintrag{Spekulationsfunktion}{Aufgabe des Lagers, Einkauf zu günstigem Zeitpunkt zu erlauben (z.\,B.\ vor Preiserhöhungen).}

\glossareintrag{Eigenlager}{Im Eigentum des Unternehmens stehendes Lager --- volle Kontrolle, hohe Fixkosten.}

\glossareintrag{Fremdlager}{Angemietetes Lager bei einem Logistikdienstleister --- variable Kosten, geringere Investition.}

\glossareintrag{Hochregallager}{Lager mit Regalen bis zu 50~m Höhe, meist automatisiert --- maximale Raumnutzung.}

\glossareintrag{Blocklager}{Lager ohne Regale, Paletten gestapelt --- günstig, für gleichartige Waren.}

\glossareintrag{Kapitalbindungskosten}{Opportunitätskosten für Kapital, das in Lagerbeständen gebunden ist und nicht produktiv investiert werden kann.}

\glossareintrag{Lagerumschlagshäufigkeit (LUH)}{Kennzahl: Wie oft pro Jahr wird der Lagerbestand komplett umgeschlagen? Formel: Jahresverbrauch $\div$ Ø~Lagerbestand.}

\glossareintrag{Durchschnittlicher Lagerbestand}{(Anfangsbestand $+$ Endbestand) $\div$ 2 oder mit 12 Monatsbeständen genauer.}

\glossareintrag{Durchschnittliche Lagerdauer}{Wie lange liegt eine Ware im Schnitt im Lager? Formel: 360 Tage $\div$ LUH.}

\glossareintrag{Lagerzinssatz}{Prozentualer Zinsanspruch auf das im Lager gebundene Kapital --- spiegelt die entgangenen Zinsen aus alternativen Investitionen.}

\glossareintrag{Lagerzinsen}{Konkrete Geldgröße: Ø~Lagerwert $\times$ Lagerzinssatz $=$ jährliche Kapitalbindungskosten.}

\end{document}
